% ============================================================
%  Chapter 7: Discussion and Synthesis
%  Henry Kabuye PhD Thesis — Northumbria University
% ============================================================

\chapter{Discussion and Synthesis}
\label{chap:discussion}

\lettrine{T}{his} chapter synthesises the findings of the three empirical investigations
presented in Chapters~3, 4, and~5, situating them within the broader literature reviewed in
Chapter~2 and addressing the research questions stated in Chapter~1.
Where individual empirical chapters reported results in isolation, this chapter analyses them
comparatively — examining how performance, uncertainty behaviour, and operational
characteristics evolve as the framework progresses from a single uncertainty-aware agent,
through a two-agent triage system, to the full three-agent knowledge-augmented architecture.
The chapter also addresses the limitations of the research programme as a whole, distinguishes
contributions to theoretical knowledge from contributions to engineering practice, and
identifies the conditions under which the proposed frameworks would and would not be expected
to transfer to operational deployment.

% ============================================================
\section{Evolution of the Detection Framework}
\label{sec:disc_evolution}
% ============================================================

% TODO: Compare detection performance across all three chapters on shared evaluation axes.
% Key points to address:
%  - How accuracy, AUC, and balanced accuracy change from Chapter 3 → 4 → 5
%  - How the escalation rate evolves as the triage policy becomes more sophisticated
%  - How calibration quality (ECE) changes as temperature scaling and peer consultation
%    are added
%  - Whether the addition of each new component (Knowledge Agent, Peer Consultation,
%    Expert Feedback Loop) produces a measurable performance gain over the preceding system

% ============================================================
\section{Uncertainty as an Executive Control Variable}
\label{sec:disc_uncertainty}
% ============================================================

% TODO: Cross-cutting analysis of the role of uncertainty across all three studies.
% Key points to address:
%  - In Chapter 3, uncertainty is a diagnostic signal reported alongside accuracy
%  - In Chapter 4, uncertainty is operationalised as a triage threshold — a direct
%    controller of autonomous versus human-reviewed decisions
%  - In Chapter 5, uncertainty drives inter-agent communication (Peer Consultation)
%    and indirectly drives the Expert Feedback Loop (escalated cases are uncertain cases)
%  - Discuss whether this progressive operationalisation of uncertainty is a general
%    principle applicable beyond ransomware detection

% ============================================================
\section{Knowledge Integration and Neuro-Symbolic Design}
\label{sec:disc_neurosym}
% ============================================================

% TODO: Discuss the Knowledge Agent's contribution relative to the two data-driven agents.
% Key points to address:
%  - How much does the Knowledge Agent improve fusion performance on the held-out test set?
%  - Under what conditions does it help most (novel families? low-confidence semantic?)
%  - Limitations of the RansKnow dataset (YouTube-sourced, 307 samples, labelling heuristic)
%  - Relationship to the neuro-symbolic AI literature (Besold et al. 2021)
%  - Future directions for richer knowledge representation (knowledge graphs, LLM-derived)

% ============================================================
\section{Continual Learning and Operational Realism}
\label{sec:disc_continual}
% ============================================================

% TODO: Discuss the Expert Feedback Loop's behaviour under simulated replay.
% Key points to address:
%  - How much does the feedback loop reduce calibration error and improve balanced accuracy
%    on held-out test sets compared to the frozen Decision MLP?
%  - Limitations of the simulated replay evaluation vs live deployment
%  - The assumption that analyst labels are accurate — implications of label noise
%  - Relationship to the concept drift literature (Ceschin 2022, Fernando 2024)

% ============================================================
\section{Research Gaps Revisited}
\label{sec:disc_gaps}
% ============================================================

% TODO: Return to the seven research gaps identified in Chapter 2 and assess the degree
% to which each has been addressed.
% Structure as: Gap → What was claimed → What the results show → Residual limitations
% Use a summary table (Gap, Status, Remaining Open Questions)

% ============================================================
\section{Limitations of the Research Programme}
\label{sec:disc_limitations}
% ============================================================

% TODO: Honest discussion of the programme-level limitations. Key items:
%  - All evaluation is on static benchmark datasets; no live deployment evaluation
%  - RanSAP and RDset represent a snapshot of the ransomware landscape; family coverage
%    may not generalise to families not represented in the benchmarks
%  - RansKnow is small (307 samples) and labelled by a heuristic threshold rule
%  - The framework has not been evaluated under adversarial attack (evasion attempts)
%  - Inference latency figures are from a research environment; production latency
%    would depend on hardware and integration architecture
%  - No formal user study to validate the usefulness of counterfactual explanations
%    with real analysts

% ============================================================
\section{Implications for Theory and Practice}
\label{sec:disc_implications}
% ============================================================

% TODO: Discuss the theoretical and practical implications of the findings.
% Theoretical:
%  - The operationalisation of epistemic uncertainty as an executive variable is a
%    general design principle applicable beyond ransomware detection
%  - Score-space CTGAN augmentation is a viable strategy for any multi-agent system
%    with low-dimensional fused representations
%  - The Peer Consultation Protocol generalises to any multi-agent system where agents
%    have heterogeneous confidence profiles under distributional shift
% Practical:
%  - Organisations seeking to deploy automated ransomware containment should treat the
%    escalation rate as a first-class design parameter alongside detection accuracy
%  - The low parameter count of the Decision MLP (~4,000 parameters) makes the Expert
%    Feedback Loop practical for deployment with modest computational resources

% ============================================================
\section{Summary}
\label{sec:disc_summary}
% ============================================================

% TODO: One-paragraph synthesis pulling together the key insights from this chapter.
% Connect back to research questions and point forward to Conclusions.
